\documentclass[11pt,a4paper]{article}

\usepackage[margin=1in]{geometry}
\usepackage[utf8]{inputenc}
\usepackage[T1]{fontenc}
\usepackage{lmodern}
\usepackage{microtype}
\usepackage{hyperref}
\usepackage{amsmath,amssymb,amsfonts,amsthm}
\usepackage{mathtools}
\usepackage[nameinlink,noabbrev]{cleveref}
\usepackage{booktabs}
\usepackage{array}
\usepackage{longtable}
\usepackage{enumitem}
\usepackage{siunitx}
\usepackage{listings}
\usepackage{xcolor}
\usepackage{caption}
\usepackage{subcaption}
\usepackage{float}
\usepackage{parskip}
\usepackage{tikz}
\usetikzlibrary{shapes.geometric, arrows, positioning}

\hypersetup{
    colorlinks=true,
    linkcolor=blue,
    urlcolor=cyan,
    citecolor=blue,
    pdftitle={InterLink: A Zero-Knowledge Atomic Cross-Chain Interoperability Protocol},
    pdfauthor={MeridianAlgo},
    pdfsubject={Blockchain Interoperability Research Paper},
    pdflang={en}
}

\definecolor{codegreen}{rgb}{0,0.6,0}
\definecolor{codegray}{rgb}{0.5,0.5,0.5}
\definecolor{codepurple}{rgb}{0.58,0,0.82}
\definecolor{backcolour}{rgb}{0.95,0.95,0.92}

% Define Rust language for listings
\lstdefinelanguage{Rust}{
  morekeywords={abstract,as,async,await,become,box,break,const,continue,crate,do,dyn,else,enum,extern,false,fn,for,if,impl,in,let,loop,macro,match,mod,move,mut,pub,ref,return,self,Self,static,struct,super,trait,true,type,typeof,unsafe,unsized,use,virtual,where,while,yield},
  sensitive=true,
  morecomment=[l]{//},
  morecomment=[s]{/*}{*/},
  morestring=[b]",
  morestring=[b]',
}
\lstdefinestyle{mystyle}{
    backgroundcolor=\color{backcolour},
    commentstyle=\color{codegreen},
    keywordstyle=\color{magenta},
    numberstyle=\tiny\color{codegray},
    stringstyle=\color{codepurple},
    basicstyle=\ttfamily\footnotesize,
    breaklines=true,
    captionpos=b,
    keepspaces=true,
    numbers=left,
    numbersep=5pt,
    showstringspaces=false,
    tabsize=2,
    language=Rust,
    frame=single,
    framesep=5pt,
    framerule=0.5pt
}
\lstset{style=mystyle}

\setlength{\parindent}{0pt}
\setlength{\parskip}{8pt plus 2pt minus 1pt}

\title{
    \vspace{2in}
    \textbf{\Huge InterLink Protocol}\\
    \vspace{0.5in}
    \large The Missing Link: A Zero-Knowledge Framework for Unified Liquidity\\
    \vspace{0.2in}
    \textit{Trustless. Scalable. Universal.}\\
    \vspace{1in}
    \textbf{MeridianAlgo Research Lab}\\
    \small Version 7.0 --- Technical Compendium --- February 2026
}
\author{}
\date{}

\begin{document}

\maketitle
\begin{titlepage}
\maketitle
\end{titlepage}
\newpage

\tableofcontents
\newpage

\phantomsection
\addcontentsline{toc}{section}{Abstract}
\begin{abstract}
\section*{Abstract}
The blockchain industry stands at a crossroads. While the proliferation of Layer-1 and Layer-2 networks has driven innovation, it has also fractured the global liquidity pool into isolated silos. Today, a user's capital on Ethereum is effectively useless on Solana, and vice versa, without traversing complex, high-latency, and often insecure bridges. The prevailing "multi-chain" reality is marred by a reliance on centralized intermediaries—federated multisigs that have become the single largest point of failure in the crypto-economy, accounting for over \$2.5 billion in losses since 2021.

\textbf{InterLink (ILINK)} proposes a paradigm shift: from trust-based interoperability to \textbf{math-based verification}. By leveraging recursive \textbf{zk-SNARKs} (specifically the Halo2 proving system), InterLink allows blockchains to cryptographically verify each other's state without human intervention. We introduce a novel "Hub-and-Spoke" architecture where a high-performance Solana Execution Hub acts as the global clearinghouse for cross-chain intent, settling transactions atomically across heterogeneous environments.

This paper details the cryptographic engineering, economic game theory, and governance structures that underpin InterLink. We present a system that not only solves the security crisis of bridging but also unlocks a new class of "Omnichain" applications—enabling developers to build software that lives everywhere at once. InterLink is not just a bridge; it is the TCP/IP layer for the decentralized web.

\vspace{0.5in}
\textbf{Keywords:} Zero-Knowledge Proofs, Recursive SNARKs, Atomic Composability, Rust, Solana Anchor, Unified Liquidity, Tokenomics, Game Theory, Interoperability, DeFi.
\end{abstract}

\newpage

\section{Introduction}
\label{sec:intro}

We are witnessing the Cambrian explosion of decentralized networks. From the Turing-complete flexibility of Ethereum to the parallel processing power of Solana and the modular sovereignty of Cosmos, developers have more choices than ever. However, this diversity comes at a steep cost: \textbf{Fragmentation}.

\subsection{The Digital Archipelago}
Imagine a world where the internet worked like blockchains do today. An email sent from a Google server would be unreadable by a Microsoft server. To send a message, you would need to pay a third-party courier to print it out, drive it to the other data center, and type it in manually. This is the state of cross-chain interoperability.

\begin{itemize}
    \item \textbf{Capital Inefficiency:} Liquidity is fractured. A lending protocol on Arbitrum cannot access collateral deposited on Optimism. This leads to high slippage and poor capital utilization.
    \item \textbf{User Friction:} Users are forced to manage multiple wallets, gas tokens, and RPC endpoints. The cognitive load is a barrier to mass adoption.
    \item \textbf{Systemic Risk:} To connect these chains, the industry built "bridges." But most bridges are merely multisig wallets controlled by a handful of signers. If these keys are compromised, the bridge collapses.
\end{itemize}

\subsection{The InterLink Vision}
We believe the future is not "multi-chain" (many disconnected chains) but "cross-chain native" (one unified network). \textbf{InterLink} is the infrastructure that makes this possible.

We replace the "trusted courier" with a \textbf{Zero-Knowledge Proof}. Instead of trusting a federation of signers, users trust a cryptographic proof that certifies a transaction occurred. If the math checks out, the transaction is valid. No human intervention required.

\section{Market Analysis: The Case for ZK}
\label{sec:market}

To understand why InterLink is necessary, we must analyze the failures of the current landscape.

\subsection{The History of Bridge Failures}
The "Bridge Trilemma" suggests you can only have two of three: Trustlessness, Generalizability, and Extensibility. Most incumbents sacrificed Trustlessness for speed.

\begin{itemize}
    \item \textbf{The Poly Network Hack (\$610M):} A privileged access management issue allowed a hacker to instruct the bridge to move funds.
    \item \textbf{The Ronin Bridge Hack (\$625M):} Five out of nine validator keys were compromised via social engineering. The bridge logic was sound, but the human element failed.
    \item \textbf{The Nomad Hack (\$190M):} A smart contract bug allowed users to spoof messages. Because the bridge relied on an "optimistic" window rather than immediate verification, the funds were drained before anyone could stop it.
\end{itemize}

These events prove that \textbf{human-operated bridges are unsustainable}. As the value secured by blockchains grows into the trillions, the incentive to attack these centralized points of failure becomes infinite.

\subsection{The ZK Advantage}
Zero-Knowledge technology offers the only viable path forward.
\begin{enumerate}
    \item \textbf{Trust Minimization:} You don't trust the Relayer; you verify their proof. A malicious Relayer cannot forge a valid ZK proof of an invalid state transition.
    \item \textbf{Capital Efficiency:} Unlike "Optimistic" bridges that require a 7-day challenge period (locking up capital), ZK proofs offer finality as soon as they are verified on-chain (minutes).
    \item \textbf{Cost Scaling:} With recursive proofs, the cost of verifying a batch of 1,000 transactions is roughly the same as verifying one. This unlocks micro-transactions previously impossible on mainnet Ethereum.
\end{enumerate}

\section{Core Definitions}
\label{sec:definitions}

To ensure clarity, we define the key actors and components of the InterLink ecosystem.

\begin{description}
    \item[The Hub (Solana):] The central coordination layer. We chose Solana for its high throughput (65k TPS) and low latency (400ms block times), making it the ideal "global order book" for cross-chain state.
    \item[The Spoke (External Chain):] Any connected blockchain (Ethereum, Arbitrum, Sui, etc.). Spokes are passive; they only react to authenticated messages from the Hub.
    \item[Gateway Contract:] The smart contract deployed on every Spoke. It acts as the "Embassy," holding assets in custody and validating incoming proofs.
    \item[Relayer:] An off-chain node that observes Spoke events, generates ZK proofs, and submits them to the Hub. Relayers are permissionless but must stake ILINK tokens.
    \item[zk-SNARK:] \textit{Zero-Knowledge Succinct Non-Interactive Argument of Knowledge}. A cryptographic receipt that proves a computation was performed correctly without revealing the underlying data.
    \item[Recursive Proof:] A technique where a ZK proof verifies another ZK proof. This allows us to "roll up" entire blocks of transactions into a single, tiny proof for cheap verification.
    \item[Atomic Composability:] The guarantee that a cross-chain transaction either fully succeeds or fully reverts. Funds are never left in limbo.
\end{description}

\section{Technical Architecture}
\label{sec:arch}

InterLink's architecture resembles a "Hub-and-Spoke" network topology, similar to the global air travel system.

\subsection{The Hub: Solana Execution Layer}
Why Solana? In a cross-chain system, the speed of the slowest link determines the user experience. By using Solana as the Hub, we ensure that the coordination layer adds negligible latency.
\begin{itemize}
    \item \textbf{State Registry:} The Hub maintains the "Global State Root"—a Merkle root representing the status of all connected chains.
    \item \textbf{Verification Engine:} A set of on-chain programs (smart contracts) that verify Groth16 and PLONK proofs.
    \item \textbf{Messaging Queue:} An ordered stream of outbound messages destined for Spokes.
\end{itemize}

\subsection{The Circuit Stack: Halo2}
We utilize the \textbf{Halo2} proving system, developed by the Zcash team. Halo2 is particularly powerful because it supports:
\begin{itemize}
    \item \textbf{Custom Gates:} We can write specialized logic (like Keccak-256 hashing) directly into the circuit for maximum efficiency.
    \item \textbf{Lookup Tables:} Essential for bitwise operations common in EVM compatibility.
    \item \textbf{No Trusted Setup (IPA):} For the inner layers of recursion, we use the Inner Product Argument (IPA), which requires no trusted setup ceremony, eliminating a major security risk.
\end{itemize}

\subsection{The Life of a Cross-Chain Message}
Let's trace a transaction from Ethereum to Arbitrum, routed via InterLink.

\begin{enumerate}
    \item \textbf{Initiation:} User deposits ETH into the Ethereum Gateway.
    \item \textbf{Finality Wait:} The Relayer waits for Ethereum finality (approx. 12 mins).
    \item \textbf{Proof Generation:} The Relayer generates a ZK proof attesting that "The Ethereum Consensus Layer has signed Block $N$, which contains transaction $T$."
    \item \textbf{Submission:} This proof is submitted to the Solana Hub.
    \item \textbf{Verification \& Routing:} The Hub verifies the proof. It sees the destination is Arbitrum. It generates a new "Outbound Message" event.
    \item \textbf{Delivery:} A Relayer (potentially a different one) sees the Solana event, proves it, and submits it to the Arbitrum Gateway.
    \item \textbf{Execution:} The Arbitrum Gateway releases the funds to the user.
\end{enumerate}

\section{Engineering \& Implementation Details}
\label{sec:impl}

Our engineering philosophy prioritizes safety and performance. The entire stack is written in \textbf{Rust}.

\subsection{Why Rust?}
Rust provides memory safety without garbage collection. In the context of cryptography and high-stakes financial software, this is non-negotiable. Rust's type system prevents entire classes of bugs (like null pointer dereferences and buffer overflows) that have led to massive hacks in C++ based systems.

\subsection{The Relayer Implementation}
The Relayer is a complex piece of software designed for 24/7 uptime.
\begin{lstlisting}[language=Rust, caption={Relayer Event Loop}]
async fn run_relayer_loop() -> Result<()> {
    let mut stream = eth_client.subscribe_logs(FILTER).await?;
    
    while let Some(log) = stream.next().await {
        // 1. Capture the event
        let event = parse_event(log)?;
        
        // 2. Fetch the Merkle Witness (Parallelized)
        let witness = tokio::spawn(fetch_witness(event.block_hash));
        
        // 3. Generate the Proof (CPU Intensive)
        let proof = rayon::spawn(move || prover.prove(witness));
        
        // 4. Submit to Solana
        solana_client.send_transaction(proof).await?;
    }
    Ok(())
}
\end{lstlisting}

\subsection{Optimizing Solana Verification}
Verifying a ZK proof on-chain is computationally expensive. To stay within Solana's compute budget (CUs), we employ several tricks:
\begin{itemize}
    \item \textbf{Precompiled Instructions:} We utilize Solana's native \texttt{ed25519} and \texttt{secp256k1} programs where possible.
    \item \textbf{Alt-Bn128 Compression:} We use a Groth16 wrapper for the final proof layer, reducing the on-chain verification to just 4 pairing checks, which is cheap enough for a single transaction.
\end{itemize}

\section{Tokenomics: The ILINK Engine}
\label{sec:tokenomics}

The ILINK token is designed to align incentives between users, relayers, and the protocol itself. It is not just a governance token; it is a utility token essential for network security.

\subsection{The Relayer Stake (The Stick)}
To become a Relayer, one must stake a significant amount of ILINK (e.g., \$50,000 worth). This stake acts as a bond.
\begin{itemize}
    \item If a Relayer submits a valid proof, they earn fees.
    \item If a Relayer submits an invalid proof (which the smart contract will reject), they lose a portion of their stake (Slashing).
    \item \textbf{Note:} Because of ZK, a Relayer \textit{cannot} successfully steal funds. The worst they can do is spam the network with bad proofs. Slashing makes this economically irrational.
\end{itemize}

\subsection{The Buy-Back-and-Burn (The Carrot)}
We implement a deflationary flywheel to accrue value to the token.
\begin{enumerate}
    \item \textbf{Fee Collection:} Users pay transaction fees in the source chain's native token (ETH, SOL, MATIC).
    \item \textbf{Protocol Swap:} The protocol aggregates these fees and buys ILINK on the open market.
    \item \textbf{Distribution:}
        \begin{itemize}
            \item \textbf{40\%} is burned immediately (removed from supply).
            \item \textbf{40\%} is paid to the Relayer who did the work.
            \item \textbf{20\%} goes to the DAO Treasury for future development.
        \end{itemize}
\end{enumerate}

This creates a direct link between usage and scarcity: \textbf{More usage $\rightarrow$ More buying pressure $\rightarrow$ Less supply}.

\section{Governance: The InterLink DAO}
\label{sec:governance}

InterLink is a public good, and its governance must be decentralized. The InterLink DAO holds the keys to the protocol's parameters.

\subsection{What does the DAO control?}
\begin{itemize}
    \item \textbf{Supported Chains:} Adding a new Spoke (e.g., connecting a new L2) requires a governance vote to deploy the Gateway contracts.
    \item \textbf{Fee Parameters:} Adjusting the base fee and the burn rate.
    \item \textbf{Emergency Pause:} In the event of a catastrophic bug, the DAO (via a security council) can pause the contracts.
\end{itemize}

\subsection{Voting Mechanics}
We utilize \textbf{Quadratic Voting} to prevent "whale" dominance. The cost of a vote is the square of the number of votes cast. This ensures that smaller holders have a voice and that passionate minorities can influence decisions that matter to them.

\section{Developer Experience (DevEx)}
\label{sec:devex}

A protocol is only as good as the apps built on top of it. We have obsessed over making InterLink easy to use.

\subsection{The InterLink SDK}
Developers don't need to know Rust or ZK math. They just install our npm package.

\begin{lstlisting}[language=JavaScript, caption={Sending a Cross-Chain Message}]
import { InterLink } from '@interlink/sdk';

const client = new InterLink({ signer: myWallet });

// Send a message from Ethereum to Solana
await client.send({
    fromChain: 'ethereum',
    toChain: 'solana',
    token: 'USDC',
    amount: '100.00',
    recipient: 'Solana_Wallet_Address',
    payload: { action: 'BUY_NFT', id: 420 } // Optional arbitrary data!
});
\end{lstlisting}

\subsection{Use Cases}
\begin{itemize}
    \item \textbf{Cross-Chain NFT Mints:} A project can launch an NFT collection on Ethereum but allow users to pay with SOL or MATIC.
    \item \textbf{Unified Yield Aggregators:} A DeFi protocol can automatically move user funds to whichever chain offers the highest APY, transparently in the background.
    \item \textbf{Global Identity:} A user can update their profile picture on one chain and have it propagate to all others.
\end{itemize}

\section{Roadmap}
\label{sec:roadmap}

We are executing on a multi-year vision.

\begin{itemize}
    \item \textbf{Phase 1: Foundation (Q2 2026)}
        \begin{itemize}
            \item Mainnet launch connecting Ethereum and Solana.
            \item Token Generation Event (TGE).
            \item Open-source release of the Relayer node.
        \end{itemize}
    \item \textbf{Phase 2: Expansion (Q4 2026)}
        \begin{itemize}
            \item Integration with Arbitrum, Optimism, and Base.
            \item Launch of the "Liquidity Layer" (internal AMM on the Hub).
            \item Developer Grants Program.
        \end{itemize}
    \item \textbf{Phase 3: Ubiquity (2027)}
        \begin{itemize}
            \item Cosmos IBC integration.
            \item Privacy-preserving transfers (using Aztec integration).
            \item Hardware acceleration for Relayers (FPGA partnerships).
        \end{itemize}
\end{itemize}

\section{Conclusion}

We built InterLink because we were tired of being afraid. Afraid that our bridge would get hacked. Afraid that our assets were stuck. Afraid that the decentralized future we were promised was turning into a fragmented mess.

InterLink is our answer. It is a return to first principles: \textbf{Don't Trust, Verify.}

By combining the cryptographic certainty of zero-knowledge proofs with the economic security of a robust token model, we are building a bridge that cannot fall. We invite you to join us in building the unified internet of value.

\newpage

\begin{thebibliography}{99}

\bibitem{halo2}
Electric Coin Co. \textit{The Halo2 Proving System}. 2024. \url{https://zcash.github.io/halo2/}

\bibitem{plonk}
Gabizon, A., Williamson, Z., Ciobotaru, O. \textit{PLONK: Permutations over Lagrange-bases for Oecumenical Noninteractive arguments of Knowledge}. 2019.

\bibitem{solana}
Yakovenko, A. \textit{Solana: A new architecture for a high performance blockchain}. 2018.

\bibitem{layerzero}
Zarick, R., et al. \textit{LayerZero: Trustless Omnichain Interoperability Protocol}. 2021.

\bibitem{kzg}
Kate, A., Zaverucha, G.M., Goldberg, I. \textit{Constant-Size Commitments to Polynomials and Their Applications}. 2010.

\bibitem{wormhole}
Wormhole Team. \textit{Wormhole: Generic Message Passing}. 2020.

\bibitem{nomad}
Paradigm Research. \textit{The Nomad Bridge Hack: A Post-Mortem}. 2022.

\end{thebibliography}

\end{document}
